% !TeX root = ../sections/03_solutions.tex
\subsection{対策-フーリエ級数}
\begin{exercise}[フーリエ級数]
	周期 \(2\pi\) をもち，区間 \([-\pi,\pi)\) 上で
	\[
	f(x)=x
	\]
	で与えられる周期関数 \(f(x)\) のフーリエ級数を求めよ．
	
	ただし，解答は \(5x\) の項まで書き下せ．
\end{exercise}

\begin{background}
	周期 \(2\pi\) の関数 \(f(x)\) のフーリエ級数は
	\[
	f(x)\sim \frac{a_0}{2}
	+\sum_{n=1}^{\infty}
	\left(a_n\cos nx+b_n\sin nx\right)
	\]
	の形で表される。
	
	フーリエ係数は次で定義される。
	\[
	a_0=\frac{1}{\pi}\int_{-\pi}^{\pi}f(x)\,dx,
	\]
	\[
	a_n=\frac{1}{\pi}\int_{-\pi}^{\pi}f(x)\cos nx\,dx,
	\]
	\[
	b_n=\frac{1}{\pi}\int_{-\pi}^{\pi}f(x)\sin nx\,dx.
	\]
	
	偶奇性を使うと計算が大幅に楽になる。
	関数 \(f(x)=x\) は奇関数である。
	また，
	\[
	\cos nx
	\]
	は偶関数，
	\[
	\sin nx
	\]
	は奇関数である。
	
	したがって，
	\[
	x\cos nx
	\]
	は奇関数であり，
	\[
	x\sin nx
	\]
	は偶関数である。
	
	よって，
	\[
	a_0=0,\qquad a_n=0
	\]
	となり，計算すべき係数は \(b_n\) だけである。
\end{background}
\begin{strategy}
	\(f(x)=x\) の場合，
	\[
	b_n=\frac{1}{\pi}\int_{-\pi}^{\pi}x\sin nx\,dx
	\]
	を計算すればよい。
	
	被積分関数 \(x\sin nx\) は偶関数なので，
	\[
	b_n=\frac{2}{\pi}\int_{0}^{\pi}x\sin nx\,dx
	\]
	となる。
	
	部分積分により，
	\[
	\int_{0}^{\pi}x\sin nx\,dx
	=
	\left[-\frac{x\cos nx}{n}\right]_{0}^{\pi}
	+\frac{1}{n}\int_{0}^{\pi}\cos nx\,dx.
	\]
	
	ここで，
	\[
	\int_{0}^{\pi}\cos nx\,dx
	=
	\left[\frac{\sin nx}{n}\right]_{0}^{\pi}=0
	\]
	であるから，
	\[
	\int_{0}^{\pi}x\sin nx\,dx
	=
	-\frac{\pi\cos n\pi}{n}.
	\]
	
	\[
	\cos n\pi=(-1)^n
	\]
	より，
	\[
	b_n
	=
	\frac{2}{\pi}
	\left(
	-\frac{\pi(-1)^n}{n}
	\right)
	=
	\frac{2(-1)^{n+1}}{n}.
	\]
	
	したがって，
	\[
	f(x)=x
	\]
	のフーリエ級数は
	\[
	x\sim
	2\sum_{n=1}^{\infty}
	\frac{(-1)^{n+1}}{n}\sin nx
	\]
	である。
	
	\(5x\) の項まで書くと，
	\[
	x\sim
	2\left(
	\sin x
	-\frac{1}{2}\sin 2x
	+\frac{1}{3}\sin 3x
	-\frac{1}{4}\sin 4x
	+\frac{1}{5}\sin 5x
	+\cdots
	\right).
	\]
	
	この問題では，
	\[
	\text{奇関数} \Rightarrow \text{正弦級数だけ}
	\]
	と見抜くことが最重要である。
\end{strategy}